% Figure: the transverse shear stress distribution tau = VQ/(Ib) across
% the depth of an I-section. The stress is parabolic within each part,
% jumps upward where the width narrows from flange to web, and peaks at
% the neutral axis. Most of the shear is carried by the web.
\begin{figure}[htbp]
\centering
\begin{tikzpicture}[scale=1.0]
% I-section outline on the left
\begin{scope}[shift={(0,0)}]
  \fill[enggray!15] (-1.4,1.6) rectangle (1.4,2.0);   % top flange
  \fill[enggray!15] (-0.25,-1.6) rectangle (0.25,1.6); % web
  \fill[enggray!15] (-1.4,-2.0) rectangle (1.4,-1.6);  % bottom flange
  \draw[engblack, thick] (-1.4,2.0) -- (1.4,2.0) -- (1.4,1.6) -- (0.25,1.6)
    -- (0.25,-1.6) -- (1.4,-1.6) -- (1.4,-2.0) -- (-1.4,-2.0) -- (-1.4,-1.6)
    -- (-0.25,-1.6) -- (-0.25,1.6) -- (-1.4,1.6) -- cycle;
  \draw[enggray, dashed] (-1.7,0) -- (1.7,0);
  \node[enggray, font=\scriptsize, anchor=west] at (1.7,0) {N.A.};
\end{scope}
% shear stress profile on the right, aligned vertically with the section
\begin{scope}[shift={(4.2,0)}]
  \draw[enggray] (0,-2.0) -- (0,2.0);   % reference axis
  % flange: small, near-linear rise from tip to junction (top)
  \draw[engblue, very thick]
    plot[domain=1.6:2.0, samples=10] ({ 0.30*(2.0 - \x) }, {\x});
  % jump at flange-web junction (top): width narrows -> tau multiplies
  \draw[engblue, very thick, dotted] (0.12,1.6) -- (0.90,1.6);
  % web: parabola from junction value up to peak at N.A.
  \draw[engblue, very thick]
    plot[domain=-1.6:1.6, samples=40]
    ({ 1.55 - 0.65*(\x)*(\x)/(1.6*1.6) - 0.65 + 0.65 }, {\x});
  % bottom flange
  \draw[engblue, very thick, dotted] (0.12,-1.6) -- (0.90,-1.6);
  \draw[engblue, very thick]
    plot[domain=-2.0:-1.6, samples=10] ({ 0.30*(2.0 + \x) }, {\x});
  \node[engred, font=\scriptsize, anchor=west] at (1.6,0) {$\tau_{\max}$};
  \fill[engred] (1.55,0) circle (0.05);
  \node[engblue, font=\scriptsize, anchor=south] at (0.9,2.05) {$\tau=\dfrac{VQ}{Ib}$};
\end{scope}
\end{tikzpicture}
\caption[Transverse shear stress across an I-section]{The transverse
shear stress $\tau=VQ/(Ib)$ across the depth of an I-section. Within
each part the stress is parabolic, since $Q$, the first moment of the
area beyond the cut, grows as the cut moves toward the neutral axis. At
the flange-web junction the width $b$ drops from the flange width to the
thin web thickness, and because $\tau$ carries $1/b$ the stress jumps
upward by that width ratio, reaching its peak $\tau_{\max}$ at the
neutral axis. The flanges carry little shear and much of the bending
moment; the web carries most of the shear and little of the moment. This
division of labour is the reason a shear check on an I-beam is a web
check, and the reason the elementary $\tau=V/A_{\mathrm{web}}$
approximation, which assumes the web takes all the shear at a uniform
stress, sits close to the exact peak.}
\label{fig:vol03ch08:shear-stress-profile}
\end{figure}
